\chapter{Discussione e conclusioni}\label{ch:conclusioni}
\section{Risposta alla domanda di ricerca}
La domanda iniziale chiedeva se un modello contestuale con esplorazione potesse mantenere pertinenza aumentando novelty e diversità. Il test conferma anzitutto la necessità di una baseline forte: Greedy anticipa gli artisti futuri molto meglio di Random e Popularity. LinUCB raggiunge un NDCG@10 vicino, senza dimostrare superiorità sul riferimento personalizzato. Il passaggio dalla validation al test inverte il loro ordine numerico; il bootstrap appaiato non individua una direzione affidabile per questo piccolo scarto, né dimostra equivalenza.

Il vantaggio più netto di LinUCB riguarda la composizione delle liste. Rispetto a Greedy, propone più item mai ascoltati e più artisti nuovi, e riduce le ripetizioni dello stesso artista nella Top-10. Gli intervalli appaiati supportano la direzione di questi effetti nel campione analizzato. Il costo è una reward proxy media inferiore, con intervallo che esclude zero. La risposta alla domanda iniziale è quindi un compromesso documentato: vicinanza nell'accuratezza di ranking, maggiore apertura delle liste e minore allineamento Top-1 all'intensità degli ascolti futuri dell'artista.

L'incremento della novelty non identifica automaticamente la rilevanza dei nuovi contenuti. Un artista nuovo può non essere ascoltato nei sette giorni successivi; un artista noto può essere una proposta appropriata e desiderata. Gli esempi qualitativi rendono visibile questa asimmetria: liste diverse possono essere tutte familiari, mentre una lista ricca di artisti nuovi può avere reward nulla. Per un prodotto reale, il peso di questi aspetti richiederebbe obiettivi espliciti ed evidenze di utilizzo.

Il risultato è leggibile grazie al legame fra componenti e protocollo: stato cronologico e feedback differito definiscono l'informazione ammessa; pool comuni rendono confrontabili le opportunità; warm-up e scaler separano preparazione e valutazione; selezione in validation e congelamento preservano il ruolo del test. Macro-media e bootstrap per utente collegano infine le conclusioni all'unità di analisi scelta. Questi accorgimenti sostengono il confronto entro il protocollo adottato, senza dimostrare che ogni singola scelta architetturale sia ottimale.

\section{Limiti del protocollo e dell'architettura}
Il vincolo principale è l'assenza di esposizioni registrate. Non sono disponibili le azioni di una logging policy né le propensioni necessarie a IPS/SNIPS. A differenza di una valutazione offline su traffico randomizzato \cite{li2010offline}, l'esperimento non stima il reward controfattuale di una raccomandazione. L'ambiente assegna a ogni scelta un valore ricostruito dalla stessa sequenza di ascolti futura, che non cambia quando cambia la policy.

La reward artist-level valorizza il riascolto di qualunque traccia dell'artista selezionato. Una raccomandazione con titolo poco pertinente può così ottenere reward elevata; una buona scoperta ascoltata oltre sette giorni può ottenere zero. Il logaritmo attenua gli ascolti estremi, ma la misura rimane influenzata dall'attività e non corrisponde a soddisfazione dichiarata. La rilevanza binaria per NDCG considera invece solo la presenza dell'artista futuro: le due metriche valutano obiettivi connessi ma diversi.

L'eleggibilità seleziona utenti attivi in tutti i periodi e restringe la generalizzabilità ai profili persistenti. Il dataset, storico e composto da utenti Last.fm, non rappresenta necessariamente una piattaforma contemporanea o un ascoltatore occasionale. Gli ID mancanti e i fallback per nome preservano gli eventi, ma non risolvono sinonimie e omonimie: un nome non collegato al relativo MBID può restare un'entità distinta. La scelta di non unificare con statistiche future evita invece una riconciliazione retrospettiva nell'osservazione della policy.

La generazione lessicografica dei pool Personal e Discovery limita quali item entrano nell'azione. L'esplorazione non raggiunge contenuti non ancora osservati globalmente e il modello non usa caratteristiche audio o semantiche. Anche la ILD distingue soltanto identità artista: due artisti simili possono essere valutati come pienamente diversi. La coverage dipende dal candidato esposto, non da un catalogo universale.

Infine, il modello lineare condiviso assume che le feature bastino a spiegare una parte stabile dell'utilità. Le preferenze possono essere non lineari e variare nel tempo. Il sistema non include una procedura di forgetting o una prova di calibrazione dell'incertezza. I confronti non isolano ogni feature o il solo meccanismo esplorativo; la conclusione riguarda la pipeline implementata.

\figura{tradeoff.pdf}{.95}{Trade-off held-out fra NDCG@10, artist novelty e ILD. I punti sono aggregati macro per utente; le distanze nel grafico non costituiscono test di significatività.}{fig:tradeoff}

\section{Sviluppi futuri}
Un'estensione prioritaria sarebbe acquisire esposizioni, feedback espliciti e propensioni della logging policy, con un protocollo di raccolta che consenta valutazione controfattuale e sperimentazione online. Una prova controllata potrebbe misurare se nuove proposte producono ascolti successivi e come incidano sulla soddisfazione, senza inferirlo dal presente proxy.

Sul piano modellistico, feature musicali o embedding potrebbero rappresentare similarità oltre l'identità artista; un confronto predefinito con modelli di diversa capacità chiarirebbe il contributo della linearità. Un'ablation delle feature e dei pool, mantenendo gli stessi vincoli temporali, aiuterebbe a separare effetto del candidato, rappresentazione e bonus esplorativo. Una diversa generazione di candidati richiederebbe però un nuovo esperimento, senza reinterpretare i risultati congelati.

Sul piano statistico, più seed e periodi di test pianificati consentirebbero di studiare la stabilità temporale. L'ottimizzazione delle risorse dovrebbe continuare tramite riuso di artefatti immutabili e replay indipendenti, verificando equivalenza prima di nuove run. Il contributo del progetto è anche questa disciplina sperimentale: identità coerenti, causalità verificata, selezione antecedente al test e provenienza leggibile rendono i risultati discutibili e riproducibili.
